% ============================================================
% MathFunction.tex
% Dissertation-wide mathematical notation and environments
% ============================================================

% ---------- Theorem-like environments ----------
\newtheorem{assumption}{Assumption}
\newtheorem{remark}{Remark}
\newtheorem{lemma}{Lemma}

% ---------- Basic math operators ----------
\DeclareMathOperator{\Tr}{tr}
\DeclareMathOperator{\diag}{diag}
\DeclareMathOperator{\rank}{rank}
\DeclareMathOperator*{\argmin}{arg\,min}
\DeclareMathOperator*{\argmax}{arg\,max}
\DeclareMathOperator{\Exp}{Exp}
\DeclareMathOperator{\Log}{Log}

% ---------- Sets and number systems ----------
\newcommand{\R}{\mathbb{R}}
\newcommand{\N}{\mathbb{N}}

% ---------- Text / debug helpers ----------
\newcommand{\todo}[1]{\textcolor{red}{[TODO: #1]}}
\newcommand{\rev}[1]{\textcolor{red}{#1}}

% ---------- Vector and matrix style ----------
% Use \vect{} for vectors and \mat{} for matrices.
% Internally both are bold, but semantic commands keep the source readable.
\newcommand{\vect}[1]{\mathbf{#1}}
\newcommand{\mat}[1]{\mathbf{#1}}
\newcommand{\bvec}[1]{\bm{#1}}
\newcommand{\bmatr}[1]{\bm{#1}}

% Common vectors/matrices
\newcommand{\zero}{\mathbf{0}}
\newcommand{\eye}{\mathbf{I}}

% ---------- Transpose, inverse, inverse transpose ----------
\newcommand{\T}{^{\top}}
\newcommand{\inv}{^{-1}}
\newcommand{\invT}{^{-\top}}

% ---------- Norms and weighted norms ----------
\newcommand{\norm}[1]{\left\lVert #1 \right\rVert}
\newcommand{\sqnorm}[1]{\left\lVert #1 \right\rVert^{2}}
\newcommand{\wnorm}[2]{\left\lVert #1 \right\rVert^{2}_{#2}}

% ---------- Probability ----------
\newcommand{\Prob}{p}
\newcommand{\given}{\,|\,}
\newcommand{\Gaussian}[2]{\mathcal{N}\!\left(#1,#2\right)}

% ---------- Optimization variables ----------
\newcommand{\stateSet}{\mathcal{X}}
\newcommand{\measSet}{\mathcal{Z}}
\newcommand{\state}{\vect{x}}
\newcommand{\meas}{\vect{z}}
\newcommand{\noise}{\vect{n}}
\newcommand{\res}{\vect{e}}

% Indexed forms
\newcommand{\statei}[1]{\vect{x}_{#1}}
\newcommand{\meask}[1]{\vect{z}_{#1}}
\newcommand{\noisek}[1]{\vect{n}_{#1}}
\newcommand{\resk}[1]{\vect{e}_{#1}}

% Covariance and information matrices
\newcommand{\cov}{\bmatr{\Sigma}}
\newcommand{\info}{\bmatr{\Omega}}
\newcommand{\covk}[1]{\bmatr{\Sigma}_{#1}}
\newcommand{\infok}[1]{\bmatr{\Omega}_{#1}}

% Jacobian / Hessian
\newcommand{\Jac}{\mat{J}}
\newcommand{\Hess}{\mat{H}}

% ---------- SLAM / factor graph ----------
\newcommand{\cost}{C}
\newcommand{\prediction}{\vect{h}}
\newcommand{\predk}[1]{\vect{h}_{#1}}
\newcommand{\factor}{\phi}
\newcommand{\factorgraph}{\mathcal{G}}

% ---------- Coordinate frames ----------
\newcommand{\FrameW}{W}
\newcommand{\FrameB}{B}
\newcommand{\FrameC}{C}
\newcommand{\FrameI}{I}
\newcommand{\FrameG}{G}

% ---------- Coordinate-frame notation ----------
% Rotation matrix:
%   \Rot{A}{B} -> {}^{A}_{B}\mathbf{R}
% Meaning:
%   rotation of frame B with respect to frame A,
%   mapping coordinates from B to A:
%   \mathbf{v}^{A} = {}^{A}_{B}\mathbf{R}\mathbf{v}^{B}
\newcommand{\Rot}[2]{\prescript{#1}{#2}{\mathbf{R}}}
\newcommand{\EstRot}[2]{\ensuremath{\prescript{#1}{#2}{\hat{\mathbf{R}}}}}

% Homogeneous transformation matrix:
%   \Tf{A}{B} -> {}^{A}_{B}\mathbf{T}
\newcommand{\Tf}[2]{\prescript{#1}{#2}{\mathbf{T}}}
\newcommand{\EstTf}[2]{\ensuremath{\prescript{#1}{#2}{\hat{\mathbf{T}}}}}

% Translation vector:
%   \tvec{A}{AB} -> \mathbf{t}^{A/AB}
% Meaning:
%   vector from origin of A to origin of B, expressed in frame A.
\newcommand{\tvec}[2]{\mathbf{t}^{#1/#2}}

% Point:
%   \pt{A} -> \mathbf{p}^{A}
% Meaning:
%   point p expressed in frame A.
\newcommand{\pt}[1]{\mathbf{p}^{#1}}

% General vector expressed in a frame:
%   \vectf{v}{A} -> \mathbf{v}^{A}
\newcommand{\vectf}[2]{\mathbf{#1}^{#2}}

% General point with point name and frame:
%   \pointf{P}{A} -> \mathbf{p}_{P}^{A}
\newcommand{\pointf}[2]{\mathbf{p}_{#1}^{#2}}
% Relative measurement
\newcommand{\RelMeas}[3]{\mat{Z}^{#1}_{#2,#3}}

% ---------- Lie groups and Lie algebra ----------
\newcommand{\SO}[1]{\mathrm{SO}(#1)}
\newcommand{\SE}[1]{\mathrm{SE}(#1)}
\newcommand{\so}[1]{\mathfrak{so}(#1)}
\newcommand{\se}[1]{\mathfrak{se}(#1)}

\newcommand{\hatop}[1]{#1^{\wedge}}
\newcommand{\veeop}[1]{#1^{\vee}}
\newcommand{\skewm}[1]{\left[#1\right]_{\times}}

% Common Lie algebra vectors
\newcommand{\liephi}{\bvec{\phi}}
\newcommand{\lietheta}{\bvec{\theta}}
\newcommand{\liexi}{\bvec{\xi}}

% ---------- Plane notation ----------
% Useful for your ground-plane factor later.
\newcommand{\plane}{\bm{\pi}}
\newcommand{\normal}{\vect{n}}
\newcommand{\planed}{d}
\newcommand{\PlaneEq}{\normal\T \vect{p} + \planed = 0}
\newcommand{\PlaneDist}[1]{\normal\T #1 + \planed}

% ---------- Convenient matrix environments ----------
\newenvironment{bmat}{\begin{bmatrix}}{\end{bmatrix}}

% Equation + aligned environment
\usepackage{environ}
\NewEnviron{ee}{
\begin{equation}
    \begin{aligned}
        \BODY
    \end{aligned}
\end{equation}
}

% ---------- Reference helpers ----------
\newcommand{\figref}[1]{Fig.~\ref{#1}}
\newcommand{\tabref}[1]{Table~\ref{#1}}
